% theme: applications_of_speed
\MajorQuestionLesson{速さの利用問題(2)}
\SetRowSpace{4mm}


% 表：単位換算を含む通常問題
\TypeQuestion{連立方程式を用いて次の問いに答えなさい。}{speed_with_conversion}

\QQ{みさきさんは、6kmの河川敷コースを歩きと自転車で進んだ。1日目は30分歩いて20分自転車に乗り、2日目は60分歩いて10分自転車に乗ると、どちらもちょうどコースの終点に着いた。歩く速さと自転車の速さは、それぞれ毎時何kmか。}{\AnswerThreeLines{$x$：自転車、$y$：徒歩（km/時）}{$\begin{cases}\dfrac13x+\dfrac12y=6\\\dfrac16x+y=6\end{cases}$}{自転車：毎時12km、歩き：毎時4km}}
\QQ{調査船が48kmの航路を低速と高速で航行した。午前は低速で2時間、高速で1時間、午後は低速で60分、高速で90分航行した。低速と高速での速さは、それぞれ毎時何kmか。}{\AnswerThreeLines{$x$：高速、$y$：低速（km/時）}{$\begin{cases}x+2y=48\\\dfrac32x+y=48\end{cases}$}{高速24km/時、低速12km/時}}
\QQ{18km離れた2地点から2人が自転車で同時に向かい合って進み、36分後に出会った。一方の速さは他方より毎時10km速かった。2人の速さは、それぞれ毎時何kmか。}{\AnswerThreeLines{$x$：遅い人、$y$：速い人（km/時）}{$\begin{cases}\dfrac35x+\dfrac35y=18\\y-x=10\end{cases}$}{遅い人10km/時、速い人20km/時}}
\QQ{13.5km離れた2地点から、ゆうたさんとまいさんが向かい合って進んだ。ゆうたさんが出発して30分後にまいさんが出発し、その60分後に出会った。まいさんの速さはゆうたさんより毎分100m速かった。2人の速さは、それぞれ毎分何mか。}{\AnswerThreeLines{$x$：ゆうた、$y$：まい（m/分）}{$\begin{cases}3x+2y=450\\y-x=100\end{cases}$}{ゆうた50、まい150（m/分）}}
\QQ{旅行で、普通列車に90分、快速列車に30分乗ると150km進み、普通列車に45分、快速列車に1時間乗ると165km進んだ。普通列車と快速列車の速さは、それぞれ毎時何kmか。}{\AnswerThreeLines{$x$：快速、$y$：普通（km/時）}{$\begin{cases}\dfrac12x+\dfrac32y=150\\x+\dfrac34y=165\end{cases}$}{快速120km/時、普通60km/時}}
\QQ{自然公園で、1日目は45分歩いたあと30分自転車に乗って11km進み、2日目は30分歩いたあと45分自転車に乗って14km進んだ。歩く速さと自転車の速さは、それぞれ毎時何kmか。}{\AnswerThreeLines{$x$：徒歩、$y$：自転車（km/時）}{$\begin{cases}\dfrac34x+\dfrac12y=11\\\dfrac12x+\dfrac34y=14\end{cases}$}{歩き：毎時4km、自転車：毎時16km}}

\AnswerColumnBreak
\TypeQuestion{連立方程式を用いて次の問いに答えなさい。}{time_with_conversion}

\QQ{湖畔の道を自転車で往復した。行きは毎時12km、帰りは毎時18kmで走り、往復に50分かかった。行きと帰りにかかった時間は、それぞれ何分か。}{\AnswerThreeLines{$x$：行き、$y$：帰り（分）}{$\begin{cases}x+y=50\\12x=18y\end{cases}$}{行き：30分、帰り：20分}}
\QQ{遊覧船が同じ航路を往復した。上りは毎時6km、下りは毎時10kmで進み、往復に480分かかった。上りと下りにかかった時間は、それぞれ何時間か。}{\AnswerThreeLines{$x$：上り、$y$：下り（時間）}{$\begin{cases}x+y=8\\6x=10y\end{cases}$}{上り：5時間、下り：3時間}}
\QQ{展望台まで同じ登山道を往復した。上りは毎時3km、下りは毎時5kmで歩き、往復に128分かかった。上りと下りにかかった時間は、それぞれ何分か。}{\AnswerThreeLines{$x$：上り、$y$：下り（分）}{$\begin{cases}x+y=128\\3x=5y\end{cases}$}{上り：80分、下り：48分}}
\QQ{15km離れた公園へ向かう途中で自転車が故障した。自転車では毎時18km、その後は毎時6kmで歩き、全部で70分かかった。自転車に乗った時間と歩いた時間は、それぞれ何分か。}{\AnswerThreeLines{$x$：自転車、$y$：徒歩（分）}{$\begin{cases}x+y=70\\\dfrac{18}{60}x+\dfrac{6}{60}y=15\end{cases}$}{自転車：40分、徒歩：30分}}

\LessonPageBreak

\QQ{港から観光施設まで、はじめは毎時36kmのローカル線に乗り、途中から毎時54kmの連絡バスに乗った。39kmの道のりを50分で移動したとき、それぞれに乗った時間は何分か。}{\AnswerThreeLines{$x$：ローカル線、$y$：バス（分）}{$\begin{cases}x+y=50\\\dfrac{36}{60}x+\dfrac{54}{60}y=39\end{cases}$}{ローカル線：20分、連絡バス：30分}}
\QQ{空港の通路1400mを、はじめは毎時4.8kmで歩き、途中から毎時3kmの動く歩道に立ち止まって乗って移動した。全部で22分かかったとき、歩いた時間と動く歩道に乗った時間は、それぞれ何分か。}{\AnswerThreeLines{$x$：徒歩、$y$：動く歩道（分）}{$\begin{cases}x+y=22\\\dfrac{4.8}{60}x+\dfrac{3}{60}y=1.4\end{cases}$}{徒歩：10分、動く歩道：12分}}


% 裏上：周回問題。単位換算の有無と情報の与え方を混在させる。
\SetRowSpace{4mm}
\TypeQuestion{連立方程式を用いて次の問いに答えなさい。}{circuit}

\QQ{1周1440mの池の周りを、2人が同じ地点から同時に出発した。反対方向に歩くと8分後に初めて出会い、同じ方向に歩くと速い人が24分後に遅い人へ初めて追いついた。2人の速さは、それぞれ毎分何mか。}{\AnswerThreeLines{$x$：速い人、$y$：遅い人（m/分）}{$\begin{cases}8x+8y=1440\\24x-24y=1440\end{cases}$}{速い人120m/分、遅い人60m/分}}
\QQ{1周1.2kmの陸上トラックで2回の走行を行った。1回目は2人が同じ地点から反対方向へ同時に走り、5分後に初めて出会った。2回目は速い人が遅い人の300m後方から同じ方向へ同時に走り、5分後に追いついた。2人の速さは、それぞれ毎分何mか。}{\AnswerThreeLines{$x$：速い人、$y$：遅い人（m/分）}{$\begin{cases}5x+5y=1200\\5x-5y=300\end{cases}$}{速い人150m/分、遅い人90m/分}}
\QQ{1周600mのロボット走行コースで2台のロボットを試験した。1回目はロボットAが出発した2分後に、ロボットBが同じ地点から反対方向へ出発し、その4分後に初めて出会った。2回目は同じ地点から同じ方向へ同時に出発し、Aが10分後にBへ初めて追いついた。2台の速さは、それぞれ毎分何mか。}{\AnswerThreeLines{$x$：A、$y$：B（m/分）}{$\begin{cases}6x+4y=600\\10x-10y=600\end{cases}$}{A：毎分84m、B：毎分24m}}
\QQ{1周1.5kmの自転車コースを、2人が同じ地点から反対方向へ同時に走り、3分後に初めて出会った。出会うまでに速い人は遅い人より300m長く進んでいた。2人の速さは、それぞれ毎時何kmか。}{\AnswerThreeLines{$x$：速い人、$y$：遅い人（km/時）}{$\begin{cases}x+y=30\\x-y=6\end{cases}$}{速い人18km/時、遅い人12km/時}}

% 裏下：列車問題。通過対象と立式の型を変える。
\AnswerColumnBreak
\TypeQuestion{連立方程式を用いて次の問いに答えなさい。}{train}

\QQ{一定の速さで走る普通列車が、線路脇の標識の前を通過するのに12秒、長さ180mのホームを通過するのに30秒かかった。列車の速さは毎秒何mで、長さは何mか。}{\AnswerThreeLines{$x$：速さ（m/秒）、$y$：長さ（m）}{$\begin{cases}12x=y\\30x=y+180\end{cases}$}{速さ：毎秒10m、長さ：120m}}
\QQ{急行列車が、反対方向へ毎時36kmで走る長さ140mの貨物列車と、すれ違い始めてから完全にすれ違うまでに10秒かかった。また、長さ240mの鉄橋を渡り始めてから渡り終わるまでに20秒かかった。急行列車の速さは毎時何kmで、長さは何mか。}{\AnswerThreeLines{$x$：速さ（m/秒）、$y$：長さ（m）}{$\begin{cases}10(x+10)=y+140\\20x=y+240\end{cases}$}{速さ：毎時72km、長さ：160m}}
\QQ{一定の速さで走る貨物列車が、長さ200mのトンネルに入り始めてから完全に出るまでに20秒、長さ500mのトンネルでは40秒かかった。列車の速さは毎秒何mで、長さは何mか。}{\AnswerThreeLines{$x$：速さ（m/秒）、$y$：長さ（m）}{$\begin{cases}20x=y+200\\40x=y+500\end{cases}$}{速さ：毎秒15m、長さ：100m}}
\QQ{一定の速さで走るモノレールが、長さ90mの駅のホームを通過するのに15秒、長さ210mの高架橋を渡り始めてから渡り終わるまでに25秒かかった。モノレールの速さは毎時何kmで、長さは何mか。}{\AnswerThreeLines{$x$：速さ（m/秒）、$y$：長さ（m）}{$\begin{cases}15x=y+90\\25x=y+210\end{cases}$}{速さ：毎時43.2km、長さ：90m}}
